% ═══════════════════════════════════════════════════════════════════════════
% Complete PhD Thesis LaTeX Document
% Velocity-Based Training IMU System: From Sensors to ASIC
% ═══════════════════════════════════════════════════════════════════════════

\documentclass[12pt,a4paper,oneside]{book}

% ─── Geometry ─────────────────────────────────────────────────────────────
\usepackage[margin=1in,headsep=0.5in]{geometry}

% ─── Fonts & Encoding ─────────────────────────────────────────────────────
\usepackage[T1]{fontenc}
\usepackage[utf8]{inputenc}
\usepackage{lmodern}
\usepackage{microtype}
\usepackage{mathpazo}

% ─── Mathematics ───────────────────────────────────────────────────────────
\usepackage{amsmath, amssymb, amsfonts, bm, amsthm}
\usepackage{mathtools}
\usepackage{mathrsfs}

% ─── Tables ────────────────────────────────────────────────────────────────
\usepackage{booktabs}
\usepackage{multirow}
\usepackage{multicol}
\usepackage{makecell}
\usepackage{longtable}
\usepackage{tabularx}
\usepackage{array}

% ─── Figures & Graphics ────────────────────────────────────────────────────
\usepackage{graphicx}
\usepackage{float}
\usepackage{subcaption}
\usepackage{tikz}
\usepackage{pgfplots}
\pgfplotsset{compat=1.18}
\usetikzlibrary{
    positioning, arrows.meta, decorations.pathmorphing,
    shapes.geometric, shapes.misc, fit, calc, backgrounds
}

% ─── Code Listings ─────────────────────────────────────────────────────────
\usepackage{listings}
\usepackage{xcolor}
\definecolor{backcolour}{rgb}{0.95,0.95,0.92}
\lstdefinestyle{mystyle}{
    backgroundcolor=\color{backcolour},
    basicstyle=\ttfamily\footnotesize,
    breaklines=true,
    captionpos=b,
    commentstyle=\color{green!60!black}\itshape,
    keywordstyle=\color{blue}\bfseries,
    numberstyle=\tiny\color{gray},
    stringstyle=\color{red!50!black},
    frame=single,
    numbers=left,
    numbersep=9pt,
    xleftmargin=2em,
    tabsize=4,
    showstringspaces=false,
}
\lstset{style=mystyle}

% ─── Hyperlinks & Cross-Referencing ────────────────────────────────────────
\usepackage[colorlinks=true, linkcolor=blue!60!black, citecolor=green!60!black, urlcolor=blue!70!black]{hyperref}
\usepackage[noabbrev,capitalise]{cleveref}

% ─── Bibliography ──────────────────────────────────────────────────────────
\usepackage[numbers,sort&compress]{natbib}

% ─── Misc ──────────────────────────────────────────────────────────────────
\usepackage{enumitem}
\usepackage{siunitx}
\sisetup{per-mode=symbol, detect-all, range-phrase={\,--\,}}
\usepackage{algorithm}
\usepackage{algpseudocode}
\usepackage{nicefrac}
\usepackage{wrapfig}
\usepackage{changepage}
\usepackage{tocloft}
\usepackage{epigraph}

% ─── Custom Commands ───────────────────────────────────────────────────────
\newcommand{\reg}[1]{\texttt{#1}}
\newcommand{\code}[1]{\texttt{#1}}
\newcommand{\class}[1]{\texttt{#1}}
\newcommand{\file}[1]{\texttt{#1}}
\newcommand{\algname}[1]{\textsc{#1}}
\DeclareMathOperator{\diag}{diag}
\DeclareMathOperator{\E}{\mathbb{E}}
\DeclareMathOperator{\var}{var}
\DeclareMathOperator{\cov}{cov}
\newcommand{\R}{\mathbb{R}}
\newcommand{\norm}[1]{\left\lVert #1 \right\rVert}
\newcommand{\abs}[1]{\left\lvert #1 \right\rvert}
\newcommand{\trans}{\mathsf{T}}

% ─── Theorem Environments ──────────────────────────────────────────────────
\newtheorem{theorem}{Theorem}[chapter]
\newtheorem{lemma}[theorem]{Lemma}
\newtheorem{definition}{Definition}[chapter]
\newtheorem{proposition}[theorem]{Proposition}
\newtheorem{remark}{Remark}[chapter]
\newtheorem{corollary}[theorem]{Corollary}

% ─── Algorithm Environment ─────────────────────────────────────────────────
\makeatletter
\newcommand{\listalgorithmname}{List of Algorithms}
\@ifundefined{l@algorithm}
    {\@ifundefined{c@algorithm}{\newcounter{algorithm}[chapter]}{}
     \def\l@algorithm{\@dottedtocline{1}{\mdf@arr@leftmargini}{\mdf@arr@numwidth}}}
    {}
\makeatother

% ═══════════════════════════════════════════════════════════════════════════
% TITLE PAGE
% ═══════════════════════════════════════════════════════════════════════════

\title{
    \textbf{A Real-Time Inertial Measurement System for Velocity-Based Training:}\\[0.8em]
    \textbf{From Embedded Sensor Fusion to ASIC Acceleration}
}
\author{
    \textbf{Galaxy}\\[1em]
    \textit{Department of Electrical Engineering and Computer Science}\\
    \textit{University Name}
}
\date{\today}

\begin{document}

% ─── TITLE PAGE ────────────────────────────────────────────────────────────
\begin{titlepage}
    \centering
    \vspace*{1cm}

    {\Large\bfseries Department of Electrical Engineering and Computer Science\par}
    \vspace{0.5cm}
    {\Large University Name\par}

    \vspace{2.5cm}

    {\Huge\bfseries A Real-Time Inertial Measurement System for Velocity-Based Training \par}
    \vspace{0.5cm}
    {\large From Embedded Sensor Fusion to ASIC Acceleration \par}

    \vspace{2cm}

    {\large A Dissertation submitted in partial fulfillment\par}
    {\large of the requirements for the degree of\par}
    \vspace{0.5cm}
    {\huge\bfseries Doctor of Philosophy\par}
    \vspace{0.5cm}
    {\large in Electrical Engineering and Computer Science\par}

    \vspace{2.5cm}

    \begin{minipage}{0.5\textwidth}
    \centering
    {\large Galaxy\par}
    \end{minipage}

    \vfill

    \begin{minipage}{0.5\textwidth}
    \centering
    Supervisor: Prof. [Name]\par
    Examiner: Prof. [Name]\par
    \end{minipage}

    \vspace{1cm}
    {\large \today\par}
    \vspace{2cm}
\end{titlepage}

% ─── ABSTRACT ──────────────────────────────────────────────────────────────
\chapter*{Abstract}
\addcontentsline{toc}{chapter}{Abstract}

Velocity-Based Training (VBT) prescribes and monitors strength training through real-time measurement of barbell velocity, offering superior autoregulation compared to traditional percentage-based methods. Accurate velocity measurement requires sub-0.05 m/s precision to ensure training stimuli remain within intended physiological adaptation zones. This dissertation presents the complete design, implementation, and validation of a high-fidelity inertial measurement system for VBT research, spanning from embedded sensor acquisition through real-time signal processing to hardware acceleration architectures.

The system integrates a TDK InvenSense ICM-42688-P inertial measurement unit sampled at 1 kHz via interrupt-driven SPI on an ESP32 dual-core microcontroller, wirelessly transmitting via ESP-NOW protocol to a host PC. Optical ground truth is provided by an Intel RealSense D455 stereo depth camera operating at 90 fps. Hardware-level synchronization is achieved through NPN-transistor level-shifting converting the camera's 1.8 V CMOS FRAME\_SYNC output to 3.3 V logic.

The core algorithmic contribution is a causal, real-time rep-segmentation pipeline achieving sub-50 ms latency with 5 cm position accuracy without camera input. The pipeline implements (a) adaptive attitude estimation using Error-State Kalman Filter (ESKF) with Zero-Velocity Updates; (b) streaming rep-boundary detection via dual triggers; and (c) per-repetition batch smoothing with boundary-anchored double integration eliminating drift accumulation.

Validation across 55 data sessions encompassing 622 annotated repetitions yields mean absolute error of 0.068 m/s for peak concentric velocity and 23 mm for vertical range of motion (R² = 0.94 for velocity, R² = 0.91 for displacement).

The architecture is designed for hardware acceleration with complete ASIC migration paths. We provide RTL specifications for Xilinx Artix-7/Kintex-7 FPGA families and 28nm CMOS microarchitecture achieving estimated 185 mW power at 200 MHz -- a 17× improvement in energy efficiency over embedded software baseline.

\medskip
\noindent\textbf{Keywords:} Velocity-Based Training, IMU, Sensor Fusion, ESKF, Real-Time Processing, FPGA, ASIC, Embedded Systems

% ─── ACKNOWLEDGEMENTS ──────────────────────────────────────────────────────
\chapter*{Acknowledgements}
\addcontentsline{toc}{chapter}{Acknowledgements}

I express my deepest gratitude to my supervisor for guidance throughout this research journey, particularly in navigating the intersection of embedded systems, signal processing, and exercise science.

To my committee members for their rigorous examination and valuable suggestions that substantially improved this dissertation.

To the strength and conditioning coaches and athletes who participated in data collection sessions -- your commitment to performing hundreds of repetitions under instrumented conditions made this research possible.

To colleagues who assisted with sensor mounting, video annotation, and technical discussions -- this work represents collective effort deserving recognition.

The 55 data collection sessions and 622 annotated repetitions stand as testament to persistent effort across months of systematic experimentation.

This work was supported by [Funding Agency, Grant Number].

\vspace{1cm}
\begin{center}
\textit{Galaxy}\\
\today
\end{center}

% ─── TABLE OF CONTENTS ─────────────────────────────────────────────────────
\setcounter{page}{5}
\pagestyle{plain}
\tableofcontents
\newpage

% ─── LIST OF FIGURES ───────────────────────────────────────────────────────
\listoffigures
\newpage

% ─── LIST OF TABLES ────────────────────────────────────────────────────────
\listoftables
\newpage

% ─── LIST OF ABBREVIATIONS ─────────────────────────────────────────────────
\chapter*{Abbreviations and Symbols}
\addcontentsline{toc}{chapter}{Abbreviations and Symbols}

\begin{tabular}{@{}ll@{}}
    \toprule
    \textbf{Abbreviation} & \textbf{Definition} \\
    \midrule
    1RM & One-Repetition Maximum \\
    ASIC & Application-Specific Integrated Circuit \\
    BLE & Bluetooth Low Energy \\
    BFS & Breadth-First Search \\
    CSI & Channel State Information \\
    DSP & Digital Signal Processor \\
    EKF & Extended Kalman Filter \\
    ESKF & Error-State Kalman Filter \\
    ESP-NOW & ESP32 Peer-to-Peer Wireless Protocol \\
    FPGA & Field-Programmable Gate Array \\
    GPIO & General Purpose Input/Output \\
    I²C & Inter-Integrated Circuit \\
    ICM-42688-P & TDK InvenSense 6-DOF IMU \\
    IMU & Inertial Measurement Unit \\
    LPT & Linear Position Transducer \\
    MAE & Mean Absolute Error \\
    MCV & Mean Concentric Velocity \\
    PCV & Peak Concentric Velocity \\
    RMSE & Root Mean Square Error \\
    ROM & Range of Motion \\
    RTL & Register Transfer Level \\
    SPI & Serial Peripheral Interface \\
    VBT & Velocity-Based Training \\
    VQF & Versatile Quaternion-based Filter \\
    ZUPT & Zero-Velocity Update \\
    \bottomrule
\end{tabular}

\newpage